% Generated draft: requires mathematical and source review.
\documentclass[11pt]{article}
\usepackage[margin=25mm]{geometry}
\usepackage{fontspec}
\setmainfont{Latin Modern Roman}
\usepackage{amsmath,amssymb,mathtools,unicode-math}
\setmathfont{Latin Modern Math}
\usepackage{enumitem,needspace,xurl,hyperref}
\setlength{\parindent}{0pt}
\setlength{\parskip}{5pt}
\setlength{\emergencystretch}{3em}
\newcommand{\N}{\mathbb{N}}
\newcommand{\Z}{\mathbb{Z}}
\newcommand{\Q}{\mathbb{Q}}
\newcommand{\R}{\mathbb{R}}
\newcommand{\C}{\mathbb{C}}
\newcommand{\F}{\mathbb{F}}
\newcommand{\A}{\mathbb{A}}
\newcommand{\PP}{\mathbb{P}}
\newcommand{\E}{\mathbb{E}}
\newcommand{\eps}{\varepsilon}
\newcommand{\dd}{\,\mathrm{d}}
\begin{document}
\section*{UnsolvedMath ID: 9000005}
\textbf{Source problem code:} catalogue entry
\subsection{Definitions and mathematical statement}
Let
\[
 SL(2,\Z)=\left\{\begin{pmatrix}a&b\\c&d\end{pmatrix}:
 a,b,c,d\in\Z,\ ad-bc=1\right\}.
\]
The source refers to a function $f$ on this group and asks whether it extends naturally to $\C/SL(2,\Z)$. It does not supply the formula for $f$, the action on $\C$, the quotient topology or structure, or a map identifying the original domain with part of the proposed new domain. Thus ``extension'' is not determined by the displayed spaces alone.

In particular one must not silently replace $\C$ by the upper half-plane or by the Riemann sphere to obtain a familiar modular action. If a Möbius action is intended, its denominator and domain issues require specification; if a real-linear action on $\C\cong\R^2$ is intended, that is different quotient data. No choice is made here.

The second question is to construct analogous series for other planar lattices or tessellations. A planar lattice is a discrete rank-two subgroup of $\R^2$, but the base series, summand, indexing, normalization, convergence mode, and quantity to be recovered are absent from the record. Tessellation and analogy criteria are likewise unspecified. This retains both research directions without inventing a modular form, a summation identity, or a target constant.
\subsection{Short English statement}
Extend the source's group function to the intended modular/orbit space, and find analogous summation series for other lattices or tessellations. The function, action, and original series are essential missing definitions.

The supplied summary describes hexagonal and other tile/lattice directions as proposed analogues and reports no published resolution found. That is an OPEN-TRIAGE assessment, not an impossibility statement. The 2018 source presents this as one of five questions, with current primary-literature audit still required.
\subsection{Sources}
\begin{itemize}
\item Local UnsolvedMath registry, permanent ID 9000005, source code catalogue entry, ``Modular extension and analogous lattice series,'' proposed by Nikita Kalinin and Mikhail Shkolnikov. Dataset curation and metadata: Ulam AI and UnsolvedMath Contributors, CC BY 4.0.
\item Kalinin and Shkolnikov, \emph{The Number pi and a Summation by SL(2,Z)} (2018), Section 4, ``Questions,'' direction 5, \url{https://amj.math.stonybrook.edu/html-articles/Files-2015-2024/17-75/index.html}; HTML-text extraction, accessed 2026-07-29.
\item The supplied summary contains no separate later primary resolution or analogous-series reference. Current primary-literature status audit and redistribution terms NEEDS\_REVIEW for the publicly accessible author/publisher copy, separately from the dataset licence.
\end{itemize}
\end{document}
